% !TEX root = ../../index.tex
\chapter{Proof-of-Concept Benchmarking Framework}

\section{Introduction to MLOps and Modern Practices}

Machine Learning Operations (MLOps) extends DevOps to ML workflows, automating model training, validation, deployment, monitoring, and retraining in production while ensuring reproducibility, compliance, and cost control. Unlike traditional software, ML introduces data drift, model staleness, and resource variability, necessitating specialized tooling for end-to-end lifecycle management. Modern MLOps in SMEs emphasizes cloud-agnostic, open-source stacks leveraging GitHub Actions for CI/CD, containerization for portability, and observability for sustainability—mirroring Fortune 500 practices at minimal cost.

\subsection{Current MLOps Landscape}

Leading tools include:
\begin{itemize}
    \item \textbf{CI/CD}: GitHub Actions (free tier), GitLab CI, Jenkins—trigger builds/tests on PRs.
    \item \textbf{Orchestration}: Kubeflow, MLflow, ZenML—pipeline DAGs for experiment tracking.
    \item \textbf{Deployment}: Kubernetes (EKS/GKE/AKS), SeldonCore, KServe—inference serving.
    \item \textbf{Monitoring}: Prometheus/Grafana, Weights \& Biases—drift detection, metrics.
\end{itemize}

\subsection{Cloud Cost Management Tools}

MLOps cost overruns (GPUs \$10+/hr) demand specialized governance:
\begin{itemize}
    \item \textbf{Kubecost}: Kubernetes cost allocation per namespace/pod, forecasts \$/month [Kubernetes native].
    \item \textbf{Cast.ai}: Auto-scales/rightsizes clusters, claims 50--90\% savings via spot instances.
    \item \textbf{CloudHealth/Spot.io}: Multi-cloud GPU optimization, interruptible workloads.
    \item \textbf{Kepler + PromQL}: Custom energy-to-cost via \$/kWh rates, aligns with thesis focus.
\end{itemize}

These prevent 30--50\% waste from idle GPUs/overprovisioning, essential for SME viability.

\subsection{Structure Overview}

\section{Monorepo Architecture and Repository Structure}

Monorepos consolidate codebases under single Git repositories, enabling atomic changes, shared tooling, and consistent builds across teams—a pattern adopted by Google (entire company), Meta (1.5M files), Microsoft (Windows+Office), Uber (200+ services), and Twitter. SMEs adapt via GitHub for 10--50 use-cases without enterprise overhead.

\subsection{Monorepo Benefits (Industry Examples)}

\begin{itemize}
    \item \textbf{Google}: Single repo for 2B lines, Bazel builds Android/Chrome/Search atomically.
    \item \textbf{Meta}: Shared OSS tools across React/Instagram/PyTorch.
    \item \textbf{Uber}: 1000+ services, Nx for ML/infra consistency.
    \item \textbf{SME Fit}: GitHub Actions matrices parallelize butterfly-vision + future pipelines.
\end{itemize}

Refactors (PyTorch upgrade) touch one PR vs 10 repos; energy tooling shared across use-cases.

\subsection{PoC Structure}

Root repository: mlops-green-ai-poc

\begin{description}
\item[README.md] Pipeline setup, quickstart
\item[.github/workflows/] CI/CD YAMLs (build/test/deploy)
\item[docker/] Base Python/ML images
\item[energy-tools/] CodeCarbon/Zeus/Perun wrappers
\item[use-cases/butterfly-vision/] Image classification
\item[monitoring/] Kepler/Grafana dashboards
\item[infra/] SSH cluster scripts
\end{description}

\subsection{Per-Use-Case YAML Configuration}

Each use-case contains a \texttt{config.yaml} driving reproducibility and energy compliance:

\begin{verbatim}
model: resnet50
batch-size: 32
epochs: 50
energy-budget-kwh: 0.5
hardware: local-loq
grid-carbon-gco2-kwh: 250
\end{verbatim}

\textbf{Key fields}:
\begin{itemize}
    \item \texttt{model}: Architecture (resnet50, efficientnet)
    \item \texttt{energy-budget-kwh}: CI/CD fail threshold
    \item \texttt{hardware}: local-loq or ssh-cluster
    \item \texttt{grid-carbon-gco2-kwh}: Regional factor (DE:\@ 250g)
\end{itemize}

\subsection{Git Workflow and Agile Practices}

Branching strategy follows Git Flow for agile MLOps:

\begin{itemize}
    \item \texttt{main}: Production deployments
    \item \texttt{dev}: Staging/integration
    \item \texttt{feature/butterfly-v2}: Parallel development
\end{itemize}

PRs trigger full validation:
\begin{enumerate}
    \item Lint (black/mypy) → pass
    \item Unit tests (pytest) → 90\% coverage
    \item Energy benchmark → $<$ config budget
    \item Container build → GHCR push
    \item Deploy to dev cluster
\end{enumerate}

\subsection{Agile Scaling Benefits}

Monorepo scales effortlessly:
\begin{itemize}
    \item Add \texttt{use-cases/medical-imaging/} → auto-discovered by CI matrix
    \item Shared \texttt{docker/base} → consistent PyTorch/CUDA
    \item Energy tooling uniform across 10+ pipelines
    \item GitHub Actions parallelizes → 5min/use-case
\end{itemize}

Matches FAANG velocity (Google monorepo → 15k commits/day) at SME cost (free tier). Energy gates ensure green compliance from sprint 1.\@


\section{CI/CD Pipeline with GitHub Actions}

GitHub Actions delivers free CI/CD (2000 min/month) integrated with monorepo, automating lint/test/build/deploy across use-cases via matrix strategies—SME equivalent to GitLab Premium.

\subsection{Pipeline Stages}

Single workflow (\texttt{.github/workflows/mlops-pipeline.yml}) triggers on push/PR:\@

\begin{enumerate}
    \item \textbf{Lint/Test}: black/mypy/pytest (90\% coverage)
    \item \textbf{Energy Benchmark}: CodeCarbon-wrapped train.py vs budget
    \item \textbf{Container Build}: Docker per use-case → GHCR
    \item \textbf{Deploy}: kubectl/SSH to dev/prod clusters
\end{enumerate}

\subsection{Workflow YAML (Key Excerpts)}

\begin{verbatim}
name: MLOps Pipeline
on: [push, pull_request]
jobs:
  matrix-test:
    runs-on: ubuntu-latest
    strategy:
      matrix:
        use-case: [butterfly-vision, medical-imaging]
    steps:
    - uses: actions/checkout@v4
    - run: |
        cd use-cases/${{ matrix.use-case }}
        pip install -r requirements.txt
        black --check .
        pytest --cov=80 tests/
  energy-bench:
    needs: matrix-test
    runs-on: ubuntu-latest
    steps:
    - run: |
        cd use-cases/butterfly-vision
        python src/train.py  # Auto-logs CodeCarbon
    - uses: actions/upload-artifact@v4
      with: {name: energy-report, path: codecarbon.csv}
  deploy:
    needs: energy-bench
    if: github.ref == 'refs/heads/dev'
    runs-on: ubuntu-latest
    steps:
    - name: Deploy to Cluster
      uses: appleboy/ssh-action@v1.0.0
      with:
        host: ${{ secrets.CLUSTER_HOST }}
        key: ${{ secrets.SSH_KEY }}
        script: |
          kubectl apply -f use-cases/butterfly-vision/k8s/
\end{verbatim}

\subsection{Secrets and Triggers}

GitHub Secrets: \texttt{CLUSTER\_HOST}, \texttt{SSH\_KEY}, \texttt{GHCR\_TOKEN}.

Triggers:
\begin{itemize}
    \item PR to dev → full pipeline + staging deploy
    \item Push to main → prod deploy (manual approval)
    \item Feature branches → test/energy only
\end{itemize}

\subsection{SME Advantages}

\begin{itemize}
    \item Parallel matrix → 10 use-cases in 8min
    \item Artifacts persist energy CSVs for Grafana
    \item Free caching → 50\% faster rebuilds
    \item SSH-action enables hybrid local/cluster
\end{itemize}

Energy fail-fast prevents wasteful deploys; badges show pipeline status on README.\@

\section{Energy Measurement Tools Integration}

Three complementary tools provide hierarchical energy observability: code-level decorators (CodeCarbon/Zeus/Perun) for per-experiment tracking, aggregated by Kepler at container/pod level for MLOps dashboards.

\subsection{Code-Level Tools Comparison}

\begin{table}[htbp]
\footnotesize
\centering
\caption{Code-Level Energy Tools}
\begin{tabular}{p{2cm} p{3.5cm} p{3cm}}
\toprule
Tool & Strengths & Limitations \\
\midrule
CodeCarbon & Easy RAPL/carbon & CPU bias \\
Zeus & Epoch stats & Benchmark-only \\
Perun & Fine traces & Complex setup \\
\bottomrule
\end{tabular}
\end{table}

\textbf{Integration Strategy}:
\begin{itemize}
\item \textbf{CodeCarbon}: Default decorator on \texttt{train.py} (CSV to MLflow)
\item \textbf{Zeus}: MLPerf validation (resnet50 benchmark)
\item \textbf{Perun}: Custom perf traces
\end{itemize}

Usage example (\texttt{use-cases/butterfly-vision/src/train.py}):

\begin{verbatim}
from codecarbon import EmissionsTracker
tracker = EmissionsTracker()
tracker.start()
# PyTorch training loop
tracker.stop()
# Fails CI if > config.yaml budget
\end{verbatim}

\subsection{Container-Level: Kepler}

Kepler DaemonSet exports pod energy to Prometheus:

\begin{verbatim}
kubectl apply -f monitoring/kepler.yaml
\end{verbatim}

Dashboard query: \texttt{sum (container\_energy\{namespace=`mlops'\}) by (pod)} visualizes butterfly-vision vs future pipelines.

\subsection{CI/CD Energy Gating}

GitHub Actions sequence:
\begin{enumerate}
\item CodeCarbon: Run training, fail if $>0.5$ kWh (config.yaml)
\item Artifact: Upload CSV to GitHub for Grafana
\item Kepler: Deploy container, validate overhead vs baseline
\end{enumerate}

Multi-tool approach captures code efficiency and orchestration cost, enabling perf/kWh optimization across use-cases.\@

\section{Butterfly Dataset and ML Pipeline}

The Butterfly image classification pipeline serves as representative CV workload, enabling end-to-end MLOps validation from data ingestion through energy-monitored serving.

\subsection{Dataset Description}

Butterfly Dataset (Kaggle: 50k images, 5 classes: Monarch, Swallowtail, etc.) mimics production defect detection:
\begin{itemize}
    \item Train: 40k images (224$\times$224 JPEG)
    \item Val/Test: 5k each
    \item Augmentation: Rotation/flip for robustness
    \item Size: 4GB → fits local LoQ, scales to cluster
\end{itemize}

Pipeline: preprocess → train ResNet50 → export ONNX → FastAPI serving.

\subsection{ML Pipeline (\texttt{use-cases/butterfly-vision/src/train.py})}

\begin{enumerate}
    \item \textbf{Data}: PyTorch DataLoader with GPU prefetch, augmentation
    \item \textbf{Train}: CodeCarbon-wrapped ResNet50 (50 epochs, AdamW)
    \item \textbf{Evaluate}: Val accuracy $>85\%$, perf/kWh logged
    \item \textbf{Export}: TorchScript/ONNX artifact
    \item \textbf{Serve}: FastAPI /predict (batch 1/32)
\end{enumerate}

Hyperparameters from config.yaml; early-stop if energy budget exceeded.

\subsection{Hardware Setup}

Dual-environment deployment validates local-to-cluster portability:

\begin{itemize}
    \item \textbf{Local}: Lenovo LOQ 15IAX9 (i5--12450HX, 16GB DDR5, 1TB NVMe, integrated Iris Xe)—development/prototyping
    \item \textbf{Cluster}: Remote server via SSH (Docker/K8s, multi-GPU planned)
\end{itemize}

SSH setup (\texttt{infra/deploy.sh}):

\begin{verbatim}
#!/bin/bash
ssh user@cluster-host << EOF
  docker pull ghcr.io/user/mlops:butterfly
  docker run -p 8080:8080 mlops:butterfly
EOF
\end{verbatim}

\textbf{CI/CD Mapping}: Local tests on LoQ specs; cluster deploys via SSH-action.

\subsection{Pipeline Validation Metrics}

Success criteria:
\begin{itemize}
    \item Accuracy: $>$88\% top-1 val
    \item Energy: $<$0.5 kWh (CodeCarbon)
    \item Latency: $<$50ms inference (cluster)
    \item Overhead: Kepler pod energy $<$1.2x bare Docker
\end{itemize}

Butterfly simulates production CV serving while fitting resource constraints, benchmarking energy attribution accuracy central to this thesis.\@

\section{Containerization and Kubernetes Deployment}

Docker containers ensure reproducibility across local LoQ and SSH clusters, orchestrated via lightweight Kubernetes (k3s) with Kepler for energy observability.

\subsection{Docker Setup}

Per-use-case Dockerfile builds slim images:

\begin{lstlisting}[language=Bash, basicstyle=\ttfamily\small]
FROM python:3.11-slim
COPY requirements.txt .
RUN pip install torch torchvision fastapi uvicorn
COPY src/ /app/src
CMD ["uvicorn", "app.main:app", "--host", "0.0.0.0"]
\end{lstlisting}

CI builds/pushes to GHCR;\@ multi-stage reduces 2GB to 800MB.\@

\subsection{Kubernetes (k3s) Cluster}

Remote SSH cluster runs k3s (lightweight K8s):

\begin{lstlisting}[language=Bash, basicstyle=\ttfamily\small]
curl -sfL https://get.k3s.io | sh -
kubectl apply -f monitoring/kepler.yaml
\end{lstlisting}

Deployment YAML structure:

\begin{itemize}
\item Deployment: 2 replicas, 2CPU/4Gi limits
\item Service: Load balancer port 80→8080
\item HPA:\@ Auto-scale on CPU $>$70\%
\end{itemize}

\subsection{Kepler Energy Monitoring}

Kepler DaemonSet deployment:

\begin{enumerate}
\item \texttt{kubectl apply -f kepler/kepler.yaml}
\item RAPL/NVML metrics via eBPF/cAdvisor
\item Prometheus export, Grafana dashboard
\end{enumerate}

Core metric: \texttt{container\_energy\{pod=`butterfly-vision'\}}

\subsection{Deployment Workflow}

\begin{enumerate}
\item GitHub Actions: Docker build/push to GHCR
\item SSH-action: \texttt{kubectl apply -f k8s/}
\item Port-forward: \texttt{kubectl port-forward svc/butterfly 8080:80}
\item Monitor: Grafana kWh/pod realtime
\end{enumerate}

Validates container overhead (10--20\%) vs bare Docker, proving energy attribution accuracy.\@
